\paragraph{Dense regional fields.} Five regional panels were re-evaluated with the operative checkpoint paths in S12. All 32-bearing profiles are retained and fitted with the primary harmonic estimator. The two CSP grids use $0.1^\circ$ spacing; the other grids use $0.5^\circ$. Line segments are transformed from north--east bearings to longitude--latitude chart directions, accounting for latitude and axis scaling. Pink crosses identify the sampled-minimum fallback, and bearing ticks are omitted at those sites. We subtract each latitude row's mean log-anisotropy and average the Pearson correlations between row segments separated by each longitude lag. The first sampled lag with mean correlation below 0.5 is reported below. These are regional, longitude-direction summaries, resolved only at the stated grid spacing; they do not estimate an isotropic global range.
\begin{table}[htbp]\centering\small
\caption{Sampling and longitude-direction correlation scales of the current fitted regional fields. The half-correlation lag is in degrees of longitude.}
\label{tab:dense}
\begin{tabular}{@{}lrrrr@{}}\toprule
Region / checkpoint&Grid ($^\circ$)&Sites&Half-corr. lag&Fallbacks\\\midrule
CSP-iNat: North America&0.1&60,501&0.2&129\\
CSP-iNat: Central Africa&0.1&40,401&0.2&79\\
SINR cap-1000&0.5&11,421&1.5&175\\
Climplicit&0.5&11,421&1&0\\
SatCLIP R18-L40&0.5&11,421&2&3\\
\bottomrule\end{tabular}
\end{table}
